% BHU PYQ -- Manually transcribed & error-corrected
\documentclass[11pt,a4paper]{article}

\usepackage[a4paper,top=2.5cm,bottom=2.5cm,left=2.8cm,right=2.8cm,headheight=14pt]{geometry}
\usepackage{amsmath,amssymb}
\usepackage[T1]{fontenc}
\usepackage[utf8]{inputenc}
\usepackage{lmodern}
\usepackage{microtype}
\usepackage{fancyhdr}
\usepackage{enumitem}
\usepackage{hyperref}
\usepackage{lastpage}
\usepackage{parskip}

\hypersetup{colorlinks=true,urlcolor=black,linkcolor=black,
  pdftitle={CHB-502: Inorganic Chemistry-III -- BHU PYQ},pdfauthor={Banaras Hindu University}}

\pagestyle{fancy}\fancyhf{}
\renewcommand{\headrulewidth}{0.4pt}\renewcommand{\footrulewidth}{0.4pt}
\fancyhead[L]{\small\textit{Banaras Hindu University}}
\fancyhead[R]{\small\textit{CHB-502: Inorganic Chemistry-III}}
\fancyfoot[C]{\small Page \thepage\ of \pageref{LastPage}}

\newlist{parts}{enumerate}{2}
\setlist[parts,1]{label=(\alph*),leftmargin=*,itemsep=5pt,topsep=3pt}
\setlist[parts,2]{label=(\roman*),leftmargin=*,itemsep=3pt,topsep=2pt}
\newcommand{\pts}[1]{\hfill\textit{[#1]}}

\begin{document}

\begin{center}
  {\large\bfseries BANARAS HINDU UNIVERSITY}\\[2pt]
  {\normalsize B.Sc. (Hons.) Semester V Examination, 2022-23}\\[6pt]
  \rule{0.8\linewidth}{0.5pt}\\[6pt]
  {\large\bfseries Paper No. CHB-502: Inorganic Chemistry-III}\\[4pt]
  {\normalsize Time: 3 Hours\hfill Maximum Marks: 70}
\end{center}
\rule{\linewidth}{0.4pt}
\small
\noindent\textit{Instructions: Answer five questions in total, including Question~1 which is compulsory.}
\normalsize
\bigskip

\noindent\textbf{Question 1.} Answer all parts of the following: \pts{5 $\times$ 2 = 10}
\begin{parts}
  \item Display the effect of a catalyst in a thermodynamically favored chemical reaction by drawing energy profile diagrams.
  \item Which statement about EDTA complexes of $\text{Mg}^{2+}$ and $\text{Ca}^{2+}$ is correct? (i) EDTA complex of $\text{Ca}^{2+}$ is more stable than $\text{Mg}^{2+}$, or (ii) EDTA complex of $\text{Mg}^{2+}$ is more stable than $\text{Ca}^{2+}$. Give reasons.
  \item Arrange the following ligands in order of their increasing field strength: $\text{ox}^{2-}, \text{OH}^-, \text{H}_2\text{O}, \text{py}, \text{en}, \text{NCS}^-$.
  \item Draw and describe the structural features of $\text{Li}_4(\text{CH}_3)_4$.
  \item Draw the structure of $\text{Ti}(\eta^1\text{-C}_5\text{H}_5)_2(\eta^5\text{-C}_5\text{H}_5)_2$ and write its IUPAC name.
\end{parts}

\medskip\rule{\linewidth}{0.2pt}

\noindent\textbf{Question 2.}
\begin{parts}
  \item How is the Laporte Selection Rule relaxed in $[\text{Cr(NH}_3)_5\text{Cl}]^{2+}$ and $[\text{CoCl}_4]^{2-}$ complexes? \pts{6}
  \item Predict the empirical formula of a metal complex using the following data: $k_3/k_2 = 3/8$ and $k_4/k_3 = 9/16$, where $k_n$ denotes stepwise formation constants. \pts{8}
\end{parts}

\medskip\rule{\linewidth}{0.2pt}

\noindent\textbf{Question 3.}
\begin{parts}
  \item Explain 'trans-effect' with proper reactions. Write down the electrostatic polarization theory and $\pi$-bonding theory of the trans-effect. \pts{8}
  \item In the electronic spectrum of $[\text{Cr(NH}_3)_5\text{Cl}]^{2+}$, assign the type of transitions to the bands at $50000\,\text{cm}^{-1}$, $25000\,\text{cm}^{-1}$, and $17000\,\text{cm}^{-1}$. \pts{6}
\end{parts}

\medskip\rule{\linewidth}{0.2pt}

\noindent\textbf{Question 4.}
\begin{parts}
  \item Write the products of the following organometallic reactions: \pts{8}
  \begin{parts}
    \item $\text{R}_3\text{SnH} + \text{Me-C}\equiv\text{C-Me} \to ?$
    \item $\text{Fe(CO)}_5 + \text{C}_5\text{H}_6 \to ?$
  \end{parts}
  \item The enthalpy of hydration ($\Delta H_{\text{hyd}}$) of $\text{Cr}^{2+}$ is $-460\,\text{kcal/mol}$. In the presence of CFSE, discuss how hydration energy is affected by crystal field effects. \pts{6}
\end{parts}

\medskip\rule{\linewidth}{0.2pt}

\noindent\textbf{Question 5.}
\begin{parts}
  \item Differentiate between thermodynamic stability and kinetic stability of coordination complexes. \pts{6}
  \item In the electronic absorption spectrum of $[\text{Ti(H}_2\text{O)}_6]^{3+}$, along with the main absorption band at $20300\,\text{cm}^{-1}$, a shoulder band is also observed. Explain why. \pts{8}
\end{parts}

\vfill
\rule{\linewidth}{0.4pt}
\begin{center}\small
  \href{https://bhu-exam.vercel.app/}{Click here to visit our website}\\[4pt]
  \href{https://me-aryan.vercel.app/}{Click here to visit our contributor}\\[4pt]
  \href{https://quashan-me.vercel.app/}{Click here to visit our contributor}
\end{center}

\end{document}
